\documentclass[11pt,a4paper]{article}
\usepackage[T1]{fontenc}
\usepackage[utf8]{inputenc}
\usepackage{amsmath,amssymb}
\usepackage[margin=2.6cm]{geometry}
\usepackage{booktabs}
\usepackage{microtype}
\usepackage{xcolor}
\usepackage[colorlinks=true,linkcolor=blue!50!black,citecolor=blue!50!black,urlcolor=blue!50!black]{hyperref}

\newcommand{\ut}{u_\tau}
\newcommand{\Ret}{\mathrm{Re}_\tau}
\newcommand{\Reb}{\mathrm{Re}_b}
\newcommand{\dtp}{\Delta t^{+}}
\newcommand{\bu}{\mathbf{u}}
\newcommand{\bU}{\mathbf{U}}
\newcommand{\bx}{\mathbf{x}}
\newcommand{\dd}{\mathrm{d}}

\title{The \texttt{slChannel} algorithm:\\
semi-Lagrangian DNS of turbulent channel flow}
\author{G.~Cavallazzi}
\date{Documentation for \texttt{slChannel} v1.0.1 --- August 2026}

\begin{document}
\maketitle

\begin{abstract}
\noindent
This document is the theory manual of \texttt{slChannel}. It records, substep
by substep, the discrete equations the released solver computes, together
with the notation and the published results the scheme rests on.
Section~\ref{sec:setting} fixes the setting and the wall-unit conventions;
Section~\ref{sec:literature} collects what the semi-Lagrangian
Navier--Stokes literature establishes, with the governing equations and the
design decisions it forces; Section~\ref{sec:scheme} specifies the
implemented scheme, and \S\ref{sec:algorithm} within it writes out one full
time step exactly as coded in \texttt{solver.py} and \texttt{semilag.py}.
The manual is versioned with the code so that the two cannot drift apart.
Validation results and the analysis of the scheme's error mechanisms are
reported separately.
\end{abstract}
\section{Setting and notation}
\label{sec:setting}

Incompressible Navier--Stokes in Lagrangian form,
\begin{equation}
\frac{\dd\bu}{\dd t} = -\nabla p + \nu\nabla^2\bu + \mathbf{f},
\qquad \nabla\cdot\bu = 0,
\qquad
\frac{\dd\bx}{\dd t} = \bu(\bx,t).
\label{eq:ns}
\end{equation}
A semi-Lagrangian scheme advances the solution on a \emph{fixed} grid by
tracing, each step, the characteristic arriving at every grid node $\bx_a$
backwards to its departure point (``foot'') $\bx_d$, and interpolating the old
field there. Wall units use $\ut$ and $\nu$:
$\dtp = \Delta t\,\ut^2/\nu$, $\Delta x^{+}=\Delta x\,\ut/\nu$. The
\texttt{slChannel} operating points are the open channel at
$\Ret\approx182$ ($256^2\times100$, $\Delta x^{+}=7.7$) and the closed
KMM \cite{kmm87} channel at $\Reb=2792.8$
($\Ret=178.12$; $192^2\times256$ with $\Delta x^{+}=11.8$ and
$256^2\times256$ with $\Delta x^{+}=8.8$). The Eulerian reference scheme
(identical spatial discretization, AB2/CN IMEX) is capped at
$\mathrm{CFL}\approx0.28$, i.e.\ $\dtp\approx0.06$; the SL schemes run at
$\dtp=0.25$--$0.40$.

\section{What the literature proves}
\label{sec:literature}

\subsection{The premise: temporal over-resolution of Eulerian DNS}

Xiu \& Karniadakis \cite{xiu01} formalize why explicit advection wastes
resolution in \emph{time}. With $N\approx \mathrm{Re}^{3/4}$ nodes per
direction to resolve the Kolmogorov length, an eigenvalue scaling
$N^{\alpha}$ of the discrete advection operator gives
\begin{equation}
\frac{\Delta t_{\mathrm{CFL}}}{\tau_\eta}
\;\propto\; \mathrm{Re}^{1/2}\,\frac{1}{N^{\alpha}}
\;\propto\; \mathrm{Re}^{\,1/2-3\alpha/4},
\label{eq:overresolve}
\end{equation}
with $\alpha=1$ (Fourier) to $2$ (Chebyshev). At $\mathrm{Re}=10^4$ the CFL
step sits one to four orders of magnitude \emph{below} the Kolmogorov time:
the smallest spatial scales are marginally resolved while the temporal scales
are over-resolved by orders of magnitude. This is exactly the
\texttt{slChannel} premise; at our $\Ret=180$ operating point the physics
limit $\dtp\approx0.3$--$0.4$ (a fraction of the near-wall Kolmogorov time
$\tau_\eta^{+}\approx 2$--$3$, and of the sweeping-decorrelation time of the
smallest resolved eddies) sits $5$--$7\times$ above the Eulerian cap.

\subsection{The error structure of interpolated SL advection}

The rigorous convergence analysis of Falcone \& Ferretti \cite{falcone98}
(anticipated by McDonald and the meteorological literature
\cite{staniforth91}) gives, for backward integration of order $k$ and
interpolation of polynomial order $P$,
\begin{equation}
\boxed{\;
\frac{u^{n+1}-u_d^{\,n}}{\Delta t} = \frac{\dd u}{\dd t}
+ \mathcal{O}\!\left(\Delta t^{k}
+ \frac{\Delta x^{P+1}}{\Delta t}\right).
\;}
\label{eq:slerror}
\end{equation}
Three consequences are proved and repeatedly verified in
\cite{xiu01,xu02,xiu05}:
\begin{enumerate}
\item \textbf{Non-monotonicity in $\Delta t$.} The error \emph{decreases} as
the step grows while the interpolation term dominates (fewer remaps per unit
time), reaches a minimum, then grows as $\Delta t^{k}$. Xiu \& Karniadakis
map this competition explicitly (their Figs.~4--5: at intermediate spatial
resolution the minimum sits at $\Delta t\approx0.024$ for their cone problem;
at high resolution pure $\Delta t^2$ growth remains).
\item \textbf{Small-$\Delta t$ over-dissipation.} The $\Delta x^{P+1}/\Delta t$
term means an SL run at \emph{small} step accumulates interpolation error at
rate $1/\Delta t$ --- Boukir et al.'s Remark~7(ii) \cite{boukir97}. SL should
be run at large $\Delta t$ or not at all.
\item \textbf{Low-order interpolation is unusable.} SL with linear
interpolation is algebraically equivalent to first-order upwinding
(excessively dissipative); cubic is the classical minimum
\cite{staniforth91}; Giraldo showed that for spectral-element interpolation
of order $P\ge4$ dispersion and dissipation errors essentially vanish at any
Courant number \cite{giraldo98}.
\end{enumerate}
Stability is subtle: $L^\infty$-stability holds for monotone (low-order)
interpolation; for order $>2$ on \emph{equidistant} stencils von~Neumann
analysis shows instabilities can arise unless ``sliding'' stencils centred on
the departure point are used (what everyone, including \texttt{slChannel},
does); for nonuniform grids ``$L^2$ stability holds although a rigorous
theory is incomplete'' \cite{xiu05}. Two further practical constraints from
\cite{xiu05}: trajectories must not cross,
\begin{equation}
\Delta t \;<\; \big\| J_{\bu}^{-1} \big\|,
\qquad J_{\bu} = \partial \bu/\partial \bx
\label{eq:crossing}
\end{equation}
(in wall units at $\Ret=180$ the near-wall shear gives
$\Delta t^{+}\lesssim \mathcal{O}(1)$, comfortably satisfied); and boundary
feet ``may search for points outside the domain due to the large time step''
--- the proposed remedy in the literature is subcycling grid points near the
boundary with $\delta t<\Delta t$, acknowledged to complicate the
implementation.

\subsection{Second-order schemes: the Crank--Nicolson trap and the
stiffly-stable cure}

This is the core methodological content of \cite{xiu01} and
\cite{boukir97}, and it maps one-to-one onto the \texttt{slChannel} time
schemes.

\paragraph{Backward integration.} Both papers trace feet with the explicit
second-order midpoint rule; writing $\alpha=\bx_a-\bx_d$,
\begin{equation}
\alpha = \Delta t\,\bu\!\left(\bx_a - \tfrac{\Delta t}{2}\,
\bu(\bx_a,t^{n}),\; t^{n+\frac12}\right),
\label{eq:midpoint}
\end{equation}
iterated to convergence (the implicit midpoint variant gives ``almost
identical results''; the explicit form is preferred for velocity fields known
only numerically \cite{xiu01}). \texttt{slChannel}'s iterated-midpoint
departure solver is exactly this.

\paragraph{First order.} Backward-Euler characteristics,
$(u^{n+1}-u_d^{\,n})/\Delta t = (-\nabla p + \nu\nabla^2 u)^{n+1}$:
unconditionally stable, robust, first order --- the direct analogue of
\texttt{slChannel} \texttt{v1}.

\paragraph{The Crank--Nicolson trap.} The obvious second-order scheme,
\begin{equation}
\frac{\bu^{n+1}-\bu_d^{\,n}}{\Delta t}
= -\nabla p^{n+1} + \nu\nabla^2\!\left(\frac{\bu^{n+1}+\bu_d^{\,n}}{2}\right),
\qquad
\frac{\dd\bx}{\dd t} = \bu^{n+\frac12} = \tfrac32\bu^{n}-\tfrac12\bu^{n-1},
\label{eq:cnwrong}
\end{equation}
``seems to be second-order in time \dots{} however, our numerical experiments
show that it is in fact still first-order'' \cite{xiu01}: integrating
(\ref{eq:ns}) along the characteristic, the midpoint rule requires the
\emph{whole} right-hand side averaged between the departure point and the
arrival point,
\begin{equation}
\frac{\bu^{n+1}-\bu_d^{\,n}}{\Delta t}
= \tfrac12\left[\big(-\nabla p + \nu\nabla^2\bu\big)_d^{\,n}
              + \big(-\nabla p + \nu\nabla^2\bu\big)^{n+1}\right].
\label{eq:cnright}
\end{equation}
Equation~(\ref{eq:cnright}) \emph{is} \texttt{slChannel}'s \texttt{v2}
(explicit terms and the extrapolated pressure gradient averaged between
departure and arrival, trajectory-CN diffusion). The paper then reports that
the three-step splitting of (\ref{eq:cnright}) ``develops a long-term
instability \dots{} due to the explicit part of the pressure term \dots{} it
could be treated by additional projections'' --- documented on the Taylor
vortex at $\mathrm{Re}=10^6$ (their Figs.~9, 16), where the CN-type scheme
drifts over hundreds of time units while the stiffly-stable scheme tracks the
Eulerian reference. Independently, \texttt{slChannel} hit the same wall
twice: the \texttt{pc} trajectory predictor diverged until the predictor was
\emph{projected} (one extra Poisson solve --- literally the ``additional
projection''), and the \texttt{v2} incremental-pressure bookkeeping had to
enter \emph{before} the implicit solve to avoid degenerating to the
non-incremental scheme.

\paragraph{The stiffly-stable (BDF2) scheme.} Both \cite{boukir97} and
\cite{xiu01} arrive at the same cure --- replace the neutrally-stable
trapezoidal combination by a backward-differentiation formula along
characteristics:
\begin{equation}
\boxed{\;
\frac{\tfrac32\bu^{n+1} - 2\bu_d^{\,n} + \tfrac12\bu_d^{\,n-1}}{\Delta t}
= \big(-\nabla p + \nu\nabla^2\bu\big)^{n+1},
\;}
\label{eq:bdf2}
\end{equation}
with \emph{two independent feet}: $\bx_d^{\,n}$ traced over one step
$\Delta t$ and $\bx_d^{\,n-1}$ traced over two steps $2\Delta t$, both
starting from the arrival node. Boukir et al.\ freeze one trajectory field
$U^{*}=2\bu^{n}-\bu^{n-1}$ for both feet and prove (Remark~4(i)) that
\emph{continuing the far foot from the near foot} (``bending'') drops the
global order to one --- the feet must be independent solutions of the same
backward ODE. Xiu \& Karniadakis trace the near foot with the extrapolated
field $\bu^{n+1/2}=\tfrac32\bu^{n}-\tfrac12\bu^{n-1}$ and the far foot with
$\bu^{n}$; either choice preserves second order (Taylor analysis). The
three-step splitting used in \cite{xiu01,xu02},
\begin{align}
\frac{\hat\bu - 2\bu_d^{\,n} + \tfrac12\bu_d^{\,n-1}}{\Delta t} &= 0,
\label{eq:split1}\\
\frac{\hat{\hat\bu}-\hat\bu}{\Delta t} &= -\nabla p^{n+1},
\qquad
\nabla^2 p^{n+1} = \frac{1}{\Delta t}\nabla\cdot\hat\bu,
\qquad
\frac{\partial p}{\partial n}
= -\nu\,\mathbf{n}\cdot\big[\hat\bu/\Delta t\big.
\big.+ \nabla\times\boldsymbol\omega^{n+1}\big],
\label{eq:split2}\\
\frac{\tfrac32\bu^{n+1}-\hat{\hat\bu}}{\Delta t} &= \nu\nabla^2\bu^{n+1},
\label{eq:split3}
\end{align}
is non-incremental in the pressure with the rotational (Karniadakis--Israeli--Orszag
\cite{kio91}) boundary condition; the implicit solve and projection act with
the effective step $\Delta t_{\mathrm{eff}} = \tfrac{2}{3}\Delta t$. Boukir
et al.\ prove stability under a non-restrictive step condition and error
estimates of the type
\begin{equation}
\|\bu-\bu_h\| \;\le\; C\!\left(\Delta t^{2} + h^{m}
+ \frac{h^{m+1}}{\Delta t}\right)
\label{eq:boukirest}
\end{equation}
for degree-$m$ elements --- second order in time, with the same
interpolation-accumulation term as (\ref{eq:slerror}).
\texttt{slChannel}'s \texttt{bdf2} scheme is the Boukir variant (frozen
$U^{*}$, two independent feet, $\Delta t_{\mathrm{eff}}=2\Delta t/3$,
non-incremental or incremental pressure), verified at self-convergence ratios
$4.00$ (incremental) and $2.0$ (non-incremental, projection-splitting
limited).

\subsection{Strong vs.\ auxiliary form; interpolation choices}

Xiu, Sherwin, Dong \& Karniadakis \cite{xiu05} compare the \emph{strong} form
(backward particle tracking, as above) against the \emph{auxiliary} form ---
solving the pure advection problem
\begin{equation}
\frac{\partial\tilde\phi}{\partial\tau} + \bu\cdot\nabla\tilde\phi = 0,
\qquad \tilde\phi(\bx,t^{n})=\phi(\bx,t^{n}),
\qquad \tau\in[t^{n},t^{n+1}]
\label{eq:oifs}
\end{equation}
with an Eulerian substepper at its own CFL limit (the OIFS scheme of Maday,
Patera \& R{\o}nquist \cite{mpr90}), so that
$\phi_d^{\,n}=\tilde\phi(\bx_a,t^{n+1})$ without any interpolation. Verdict:
at intermediate steps the auxiliary form is preferable (no interpolation
error, standard parallelization), but its cost grows with the number of
substeps and it is \emph{not} CFL-free; \textbf{at large time steps only the
strong form is stable and its tracking cost is essentially independent of}
$\Delta t$. \texttt{slChannel} is a strong-form code; this is the right
branch for $\dtp\gtrsim0.25$.

For the channel geometry specifically (\cite{xiu05}, \S3.3: Fourier
$\times$ Fourier $\times$ spectral-element, i.e.\ structurally our grid),
they reject global Fourier interpolation as ``very high'' cost and use
\emph{local} interpolation in the homogeneous directions --- exactly the
tensor-product-stencil choice of \texttt{slChannel} --- comparing Lagrange
stencils of orders 3--7 against \emph{cubic Hermite} with the nodal
derivatives computed spectrally (FFT differentiation):
third-order Hermite matches seventh-order Lagrange accuracy at about half
the cost, and retains exponential convergence where fixed-order Lagrange
saturates. This is the closest published relative of \texttt{slChannel}'s
prefiltered-spline \texttt{field\_interp} path (C$^2$ B-spline in $x,y$ via
FFT symbol division $+$ Hermite evaluation in $z$).

\subsection{The one turbulence result: Xu, Xiu \& Karniadakis (2002)}

Reference \cite{xu02} is the only published SL turbulent channel DNS in this
line, and its parameters land startlingly close to ours:
\begin{center}
\begin{tabular}{lll}
\toprule
 & Xu et al.\ 2002 \cite{xu02} & \texttt{slChannel} (2026) \\
\midrule
Scheme & BDF2 characteristics, Eqs.~(\ref{eq:bdf2})--(\ref{eq:split3})
       & same skeleton (\texttt{bdf2}, Boukir feet) \\
Discretization & Fourier$\times$SE(Jacobi)$\times$Fourier, $48^3$
               & FD2 staggered, $192^2$--$768^2\times100$--$256$ \\
Interpolation & \emph{global} modal (all Fourier modes)
              & local tricubic/triquintic, Triton \\
$\Ret$ & 116 & 178--182 (exact KMM $\Reb$) \\
$\dtp$ & $0.2645$ ($10\times$ Eulerian $\Delta t$)
       & $0.25$--$0.40$ ($4$--$7\times$) \\
Statistics window & $T=20$ convective units
                  & 70 washouts ($\approx 880$ c.u., $10^4\,t^{+}$) \\
Validation & Eulerian twin only, mean $+$ rms
           & Eulerian twin $+$ KMM/MKM data, \\
 & & profiles, stresses, spectra \\
Cost & $10\times$/step $\Rightarrow$ net $1.0\times$
     & $1.3$--$2.1\times$/step $\Rightarrow$ net $2$--$3.2\times$ \\
\bottomrule
\end{tabular}
\end{center}
Their findings: (i) the SL run is \emph{stable at marginal resolution without
de-aliasing} where the Eulerian twin is unstable at any $\Delta t$ ---
interpolation supplies the missing high-$k$ dissipation; (ii) mean velocity
indistinguishable from the Eulerian reference ``despite the very large time
step''; (iii) \emph{rms intensities differ, most visibly an elevated
streamwise peak}, which they attribute to time-averaging error by comparing
$T=5$ vs $T=20$ Eulerian averages; (iv) the global interpolation made the
method $10\times$ slower per step, cancelling the $10\times$ step gain ---
they close by calling for ``better local interpolation procedures'' and
report no follow-up channel study. Two details deserve emphasis. First,
their $\dtp=0.2645$ sits \emph{just below} the two-foot stability boundary
$\dtp\approx0.30$ that \texttt{slChannel} mapped in 2026 --- they were one
courageous step-doubling away from a mysterious blow-up. Second, their
elevated streamwise rms is the same signature \texttt{slChannel} measures
with $44\times$ longer averaging ($+5.6\%$ peak $u'^{+}_{\mathrm{rms}}$ at
$\Delta x^{+}=11.8$); with converged statistics it is demonstrably \emph{not}
an averaging artifact.

\section{The \texttt{slChannel} scheme}
\label{sec:scheme}


\subsection{Scheme summary}

Staggered MAC finite differences ($x,y$ uniform periodic, $z$
tanh-stretched), \texttt{torch.float64} infrastructure with fp32-accumulate
Triton gather kernels for the interpolation. Feet by iterated midpoint
(\ref{eq:midpoint}) with trilinear (default) or tricubic trajectory sampling;
fields remapped by tensor-product Lagrange stencils (tricubic $P{=}3$ /
triquintic $P{=}5$; nonuniform-node weights in $z$ against the actual grid,
analytic inverse tanh map for stencil location).

The time scheme is Boukir BDF2 (\ref{eq:bdf2}) and nothing else: frozen
$U^{*}=2\bu^{n}-\bu^{n-1}$, two independent feet at $\Delta t$ and
$2\Delta t$, update $\hat u = (4\bar u - \bar{\bar u})/3$, $\theta{=}1$
$z$-solve and FFT projection at $\Delta t_{\mathrm{eff}}=2\Delta t/3$, with
\texttt{noninc} (production) or \texttt{inc} pressure --- velocity
self-convergence ratios $2.0$ and $4.00$ respectively. Mass flux is imposed
exactly by a uniform divergence-free shift (CaNS convention); Reynolds
stresses are accumulated about the \emph{time} mean at each component's own
staggered nodes.

The development campaigns also exercised a first-order
backward-Euler splitting (\texttt{v1}), a departure-averaged
Crank--Nicolson scheme with incremental pressure (\texttt{v2},
self-convergence $3.93$), projected predictor--corrector trajectories
(\texttt{pc}, $3.78$), and a C$^{2}$ prefiltered spline remap. All are
reported here as \emph{evidence}, and all were removed from the code once
BDF2 won the stability-versus-accuracy trade: the released solver implements
the single scheme above. An Eulerian IMEX path is retained purely as a
like-for-like reference for validation and benchmarking.

\subsection{One step in full: the discrete equations}
\label{sec:algorithm}

This subsection writes out every substep as implemented in
\texttt{solver.py}/\texttt{semilag.py}, in execution order. Notation:
$V^{n}=(u,v,w)^{n}$ is the staggered MAC velocity ($u$ at $x$-faces, $v$ at
$y$-faces, $w$ at $z$-faces, one ghost layer per side; pressure at cell
centres); $G_h$ and $D_h$ are the staggered gradient and divergence;
$L_{xy}$ and $L_z$ the horizontal and wall-normal discrete Laplacians
(second-order central, $L_z$ on the stretched metric); $\mathcal A[\phi](\bx_d)$
denotes the interpolation of a field $\phi$ at the departure points of the
component grid on which $\phi$ lives. All ghost layers are rebuilt
(\emph{BC application}) after every operation that leaves them stale.

\paragraph{(1) Trajectory velocity.} A single advecting field
$\bU^{*}$, defined per scheme:
\begin{equation}
\bU^{*} =
\begin{cases}
\tfrac32 V^{n}-\tfrac12 V^{n-1} & \text{\texttt{v1}/\texttt{v2}, \texttt{ab2}
(target $t^{n+1/2}$)},\\[2pt]
V^{n} & \text{\texttt{traj\_extrapolation: none} (diagnostic, $\mathcal O(\Delta t)$)},\\[2pt]
\tfrac12\big(V^{n}+\tilde V^{n+1}\big) & \text{\texttt{pc} (projected predictor
$\tilde V^{n+1}$, see below)},\\[2pt]
2V^{n}-V^{n-1} & \text{\texttt{bdf2} (frozen for \emph{both} feet, target $t^{n+1}$)}.
\end{cases}
\label{eq:ustar}
\end{equation}
For \texttt{pc}, the predictor is a complete first-order step (SL advection of
$V^{n}$ by itself, plus the horizontal viscous tendency, plus the
\emph{lagged} pressure gradient $-G_h p^{n-1/2}$, then a backward-Euler
$z$-solve) which is then \emph{projected} --- one extra Poisson solve. The
projection is not optional: an unprojected predictor carries an
$\mathcal O(\Delta t)$ divergence, SL trajectories through a divergent
mid-velocity are compressive, and the scheme blows up
($\ut$ $+127\%$).

\paragraph{(2) Departure points: iterated midpoint.} For every arrival node
$\bx_a$ of each component grid, with $K=2$ iterations
(\texttt{n\_traj\_iters}):
\begin{equation}
\bx_m^{(0)}=\bx_a,\qquad
\bx_m^{(s)} = \bx_a - \frac{\Delta t}{2}\,\bU^{*}\!\big(\bx_m^{(s-1)}\big),
\quad s=1,\dots,K,\qquad
\bx_d = 2\bx_m^{(K)}-\bx_a,
\label{eq:feet}
\end{equation}
i.e.\ $\bx_d = \bx_a - \Delta t\,\bU^{*}(\bx_m^{(K)})$: the explicit iterated
midpoint rule (\ref{eq:midpoint}). $\bU^{*}$ is sampled at the (staggered)
midpoints by trilinear interpolation on the fast path
(\texttt{traj\_order:~2}; its $C^{0}$ interpolant caps \emph{self}-convergence
at fixed grid at $\mathcal O(\Delta t)$ with a small $h^{2}$ coefficient ---
tricubic sampling, \texttt{traj\_order:~4}, restores clean
$\mathcal O(\Delta t^{2})$, ratio $3.93$). $x,y$ are unwrapped modulo the
period; $z$ is clamped to the wall-adjacent node range and every clamp is
counted (\texttt{clamped} diagnostic; measured $0$ over $10^{4}\,t^{+}$ in
the closed channel). For \texttt{bdf2} the same $\bU^{*}$ traces \emph{two
independent} feet from the same arrival node, (\ref{eq:feet}) with
$\Delta t$ and with $2\Delta t$ --- never the far foot continued from the
near one (Boukir Remark~4(i)).

\paragraph{(3) Remap.} The old field is evaluated at the feet by a
tensor-product Lagrange stencil of $P{+}1$ nodes per direction
($P=3$ tricubic, $P=5$ triquintic):
\begin{equation}
\mathcal A[\phi](\bx_d)
= \sum_{i,j,k} w^{x}_{i}(x_d)\, w^{y}_{j}(y_d)\, w^{z}_{k}(z_d)\;
\phi^{n}_{i j k},
\qquad
w_{i}(\xi)=\prod_{m\ne i}\frac{\xi-\xi_m}{\xi_i-\xi_m}.
\label{eq:remap}
\end{equation}
In $x,y$ the nodes are uniform and the weights close-form; in $z$ the weights
are built \emph{against the actual stretched nodes} $z_c$/$z_f$ (uniform-$\xi$
weights on the tanh image would silently lose an order: $z_c$ are face
midpoints, not the tanh-map image), the stencil is located by the analytic
inverse tanh map plus one node compare (no search), sliding stencils centred
on the foot in the interior, one-sided at the walls. The index/weight set of
(\ref{eq:remap}) is computed once per component and step and reused for every
field remapped on that grid --- this is what makes the \texttt{v2}
departure-evaluated right-hand sides in (\ref{eq:v2star}) nearly free. The
alternative $C^{2}$ path (\texttt{field\_interp: spline}) replaces
(\ref{eq:remap}) by a prefiltered cubic B-spline in $x,y$ (FFT symbol
division) and a nonuniform $C^{2}$ cubic spline in $z$ with Hermite
evaluation.

\paragraph{(4) Explicit horizontal tendency.} On $V^{n}$,
$R_{xy}=\nu L_{xy}V^{n}$ (no body force here --- the flux constraint is
substep 8). \texttt{v2} time-centres it by AB2,
$R_{xy}^{\,n+1/2}=\tfrac32 R_{xy}^{\,n}-\tfrac12 R_{xy}^{\,n-1}$; \texttt{bdf2}
extrapolates to $t^{n+1}$, $R_h=2R_{xy}^{\,n}-R_{xy}^{\,n-1}$
(\texttt{bdf2\_xy\_rhs: extrap}; \texttt{lagged} uses $R_{xy}^{\,n}$).

\paragraph{(5) Advected predictor.} Per scheme:
\begin{align}
\texttt{v1:}\quad & u^{*} = \mathcal A[V^{n}](\bx_d) + \Delta t\,R_{xy},
\label{eq:v1star}\\[4pt]
\texttt{v2:}\quad & u^{*} = \mathcal A[V^{n}](\bx_d)
+ \Delta t\Big[\tfrac12\big(\mathcal R_d + \mathcal R_a\big)
+ \tfrac12 \big(\nu L_z V^{n}\big)_d\Big],
\qquad
\mathcal R = R_{xy}^{\,n+1/2} - G_h\, p_{\mathrm{ext}},
\label{eq:v2star}\\[4pt]
\texttt{bdf2:}\quad & u^{*} =
\tfrac43\,\mathcal A[V^{n}](\bx_d^{\,\Delta t})
- \tfrac13\,\mathcal A[V^{n-1}](\bx_d^{\,2\Delta t})
+ \Delta t_{\mathrm{eff}}\big[R_h - G_h\,p_{\mathrm{ext}}\big],
\qquad
\Delta t_{\mathrm{eff}}=\tfrac{2}{3}\Delta t.
\label{eq:bdfstar}
\end{align}
In (\ref{eq:v2star}) subscripts $d/a$ mean the field evaluated at the
departure point (via the reused stencil of substep 3, with BC-consistent
ghosts so near-wall feet interpolate meaningful values) and at the arrival
node: the \emph{entire} explicit right-hand side, including the extrapolated
pressure gradient $p_{\mathrm{ext}}=2p^{n-1/2}-p^{n-3/2}$, is averaged along
the characteristic --- this is exactly what distinguishes the true
second-order scheme (\ref{eq:cnright}) from the first-order trap
(\ref{eq:cnwrong}). The $z$-viscous term enters with its departure half only;
its arrival half is the implicit solve below. In (\ref{eq:bdfstar}),
$p_{\mathrm{ext}}=2p^{n}-p^{n-1}$ (full levels) for
\texttt{bdf2\_pressure: inc} and is omitted for \texttt{noninc}; dividing
Boukir's update (\ref{eq:bdf2}) by the BDF2 leading coefficient $\tfrac32$ is
what produces the $\tfrac43/\tfrac13$ foot weights and the effective step
$\Delta t_{\mathrm{eff}}$ carried by every subsequent substep. The pressure
term must enter \emph{here}, before the implicit solve: adding it after is
algebraically identical to the non-incremental scheme (verified analytically
and numerically).

\paragraph{(6) Implicit wall-normal diffusion.} One tridiagonal solve per
horizontal position and component ($\theta$-scheme on the stretched metric,
no-slip rows at the walls):
\begin{equation}
\big(I - \theta\,\delta t\,\nu L_z\big)\,u^{**}
= u^{*} + (1-\theta)\,\delta t\,\nu L_z u^{*},
\qquad
(\theta,\delta t)=
\begin{cases}
(\tfrac12,\ \Delta t) & \texttt{v1 (CN at arrival)},\\
(1,\ \tfrac{\Delta t}{2}) & \texttt{v2 (arrival half of trajectory-CN)},\\
(1,\ \Delta t_{\mathrm{eff}}) & \texttt{bdf2}.
\end{cases}
\label{eq:zsolve}
\end{equation}

\paragraph{(7) Projection.} With $\delta t=\Delta t$
(\texttt{v1}/\texttt{v2}) or $\Delta t_{\mathrm{eff}}$ (\texttt{bdf2}),
\begin{equation}
L_h\,\phi = \frac{D_h\,\bu^{**}}{\delta t},
\qquad
\bu^{n+1} = \bu^{**} - \delta t\,G_h\,\phi,
\qquad
p^{\mathrm{new}} =
\begin{cases}
\phi & \texttt{v1}, \texttt{bdf2-noninc},\\
p_{\mathrm{ext}}+\phi & \texttt{v2}, \texttt{bdf2-inc}.
\end{cases}
\label{eq:proj}
\end{equation}
The Poisson solve diagonalizes $x,y$ by real FFTs with the second-order
modified wavenumbers $\tilde k^{2}=2(1-\cos k\Delta)/\Delta^{2}$, leaving one
tridiagonal system in $z$ per mode pair $(k_x,k_y)$ with homogeneous Neumann
ends (built into the staggered stencil --- no explicit pressure BC, hence no
rotational-BC machinery of (\ref{eq:split2}) is needed); the singular
Neumann--Neumann $(0,0)$ mode is pinned. On the incremental branches the
splitting error of (\ref{eq:proj}) is $\mathcal O(\Delta t^{2})$
(self-convergence ratios $3.93$/$4.00$); the non-incremental branch is
first-order in the projection but robust (ratio $2.0$).

\paragraph{(8) Exact bulk-flux constraint.} Last, the CaNS-convention mass
flux enforcement:
\begin{equation}
c = U_b - \langle u^{n+1}\rangle_V,
\qquad
u^{n+1} \mathrel{+}= c \ \ \text{(uniformly)},
\qquad
f = c/\Delta t,
\label{eq:flux}
\end{equation}
where $\langle\cdot\rangle_V$ is the volume average. A uniform shift is
divergence-free, so it commutes with the projection just performed; the
driving force $f$ is \emph{diagnosed} from the correction rather than steered
towards it. This replaces an integral controller whose error dynamics were an
undamped oscillator with $1/\Delta t$ gain --- under which any $\Delta t$
sweep partly measured the controller instead of the scheme.

\paragraph{Bootstrapping.} \texttt{bdf2} assumes constant $\Delta t$ (the
$2\Delta t$ foot and the $\tfrac43/\tfrac13$ weights); on the first step,
after a restart, or on any $\Delta t$ change the step degenerates to one BDF1
step (backward-Euler characteristics: $\bU^{*}=V^{n}$, single foot,
$\theta=1$ at full $\Delta t$) which seeds the $V^{n-1}$ and $R_{xy}^{\,n-1}$
histories --- one $\mathcal O(\Delta t^{2})$ local error, the same policy as
the documented AB2 restart re-bootstrap. \texttt{bdf2} configs pin
$\Delta t$ (\texttt{dt\_update\_interval:~0}).
\section*{Note on tooling}
\addcontentsline{toc}{section}{Note on tooling}
Parts of the implementation, the analysis scripts and the drafting of this
report were carried out with the assistance of Claude (Anthropic). All results
reported here were produced by the code in the repository and verified by the
author, who takes full responsibility for their content.
\begin{thebibliography}{9}
\bibitem{boukir97}
K.~Boukir, Y.~Maday, B.~M\'etivet, E.~Razafindrakoto,
A high-order characteristics/finite element method for the incompressible
Navier--Stokes equations,
\emph{Int.\ J. Numer.\ Meth.\ Fluids} \textbf{25}, 1421--1454 (1997).

\bibitem{xiu01}
D.~Xiu, G.E.~Karniadakis,
A semi-Lagrangian high-order method for Navier--Stokes equations,
\emph{J.~Comput.\ Phys.} \textbf{172}, 658--684 (2001).

\bibitem{xu02}
J.~Xu, D.~Xiu, G.E.~Karniadakis,
A semi-Lagrangian method for turbulence simulations using mixed spectral
discretizations,
\emph{J.~Sci.\ Comput.} \textbf{17}, 585--597 (2002).

\bibitem{xiu05}
D.~Xiu, S.J.~Sherwin, S.~Dong, G.E.~Karniadakis,
Strong and auxiliary forms of the semi-Lagrangian method for incompressible
flows,
\emph{J.~Sci.\ Comput.} \textbf{25}, 323--346 (2005).

\bibitem{falcone98}
M.~Falcone, R.~Ferretti,
Convergence analysis for a class of high-order semi-Lagrangian advection
schemes,
\emph{SIAM J. Numer.\ Anal.} \textbf{35}, 909--940 (1998).

\bibitem{giraldo98}
F.X.~Giraldo,
The Lagrange--Galerkin spectral element method on unstructured quadrilateral
grids,
\emph{J.~Comput.\ Phys.} \textbf{147}, 114--146 (1998).

\bibitem{staniforth91}
A.~Staniforth, J.~C\^ot\'e,
Semi-Lagrangian integration schemes for atmospheric models --- a review,
\emph{Mon.\ Wea.\ Rev.} \textbf{119}, 2206--2223 (1991).

\bibitem{kio91}
G.E.~Karniadakis, M.~Israeli, S.A.~Orszag,
High-order splitting methods for the incompressible Navier--Stokes equations,
\emph{J.~Comput.\ Phys.} \textbf{97}, 414--443 (1991).

\bibitem{mpr90}
Y.~Maday, A.T.~Patera, E.M.~R{\o}nquist,
An operator-integration-factor splitting method for time-dependent problems:
application to incompressible fluid flow,
\emph{J.~Sci.\ Comput.} \textbf{5}, 263--292 (1990).

\bibitem{kmm87}
J.~Kim, P.~Moin, R.~Moser,
Turbulence statistics in fully developed channel flow at low Reynolds number,
\emph{J.~Fluid Mech.} \textbf{177}, 133--166 (1987);
R.D.~Moser, J.~Kim, N.N.~Mansour,
DNS of turbulent channel flow up to $\mathrm{Re}_\tau=590$,
\emph{Phys.\ Fluids} \textbf{11}, 943--945 (1999).
\end{thebibliography}

\end{document}
